\documentclass[12pt,a4paper]{article}
\usepackage[utf8]{inputenc}
\usepackage[margin=2.5cm]{geometry}
\usepackage[numbers]{natbib}
\bibliographystyle{ieeetr}

\title{PE5 -- Vista previa individual\\Conclusión 4 --- Unidad 4: Validación y gestión de cambios}
\author{Nieves Sánchez, Jimmy Samuel}
\date{}

\begin{document}
\maketitle



\subsection*{Conclusión 4 --- Unidad 4: Validación y gestión de cambios}

La inspección formal detecta con alta eficacia los defectos de contenido para
los que fue diseñada, pero no está concebida para exponer huecos
estructurales en el propio formato de especificación, que solo emergen al
reauditar el documento completo bajo una lente distinta \citep{wohlin2012}.
Así lo evidencia el diagnóstico de trazabilidad realizado para el criterio
C2 sobre el ERS final de MundiPets: la inspección Fagan de la PE4 se acotó
explícitamente a la Sección 3 (Requisitos) del documento, por lo que nunca
revisó los modelos UML de la Sección 4. Como consecuencia, el equipo modeló
en el diagrama de clases una entidad \textit{VeterinaryAppointment} para
gestionar citas veterinarias sin que existiera ningún RF ni CU que la
respaldara y que además contradecía directamente dos exclusiones de
alcance ya declaradas en el propio ERS (``no reemplaza consulta
veterinaria'', ``no ofrece agenda médica completa''). Este defecto de tipo
cadena rota no fue detectado por la inspección Fagan por diseño, no por
descuido de los inspectores, sino hasta que se realizó el diagnóstico
sistemático de cobertura de trazabilidad de la PE5, al cruzar cada clase del
diagrama contra el ERS. Como resultado, la clase se eliminó de la partición
Mascotas y Salud, diferenciando el recordatorio pasivo (RF-10), que sí está
en alcance, de la reserva activa de turnos, que permanece fuera de él.
Consistente con lo registrado en la retrospectiva del equipo (Sección 9),
donde se identificó que ni siquiera la propia revisión de trazabilidad
estuvo exenta de una segunda pasada, el equipo adopta como práctica extender
la verificación cruzada más allá del alcance formal de la inspección de
contenido, cubriendo también la consistencia entre los modelos y el texto
del ERS, tal como exige la gestión de configuración de \citep{ieee29148}.

\bibliography{conclusion4_referencias}

\end{document}
